% !TEX program = xelatex
% GENERATED by tools/build.py from data/records.csv and templates/cv_long.tex.
% Edit the data or the template, not this file.
\documentclass[11pt, a4paper]{cvClass}
\geometry{left=2cm, top=1.2cm, right=2cm, bottom=1.6cm, footskip=.7cm}
% ---------------------------------------------------------------------------
% Shared layout macros for cv_long.tex, cv_short.tex and publication_list.tex
% (input by each template; lives in templates/, copied next to the .tex)
% ---------------------------------------------------------------------------
\colorlet{awesome}{awesome-red}          % the CV orange-red (#DC3522)
\definecolor{accent2}{HTML}{E3B23C}      % muted yellow, used sparingly
\setlength{\parskip}{0pt}
\urlstyle{same}

% A hanging list: year (or label) on the left, text on the right
\newlength{\cvlabelwidth}\setlength{\cvlabelwidth}{2.45cm}
\newlength{\cvitemsep}\setlength{\cvitemsep}{1.4mm}
\newenvironment{cvlist}{%
  \vspace{1.2mm}%
  \begin{list}{}{%
    \setlength{\leftmargin}{\cvlabelwidth}\setlength{\labelwidth}{\dimexpr\cvlabelwidth-3mm\relax}%
    \setlength{\labelsep}{3mm}\setlength{\itemsep}{\cvitemsep}%
    \setlength{\parsep}{0pt}\setlength{\topsep}{0pt}\setlength{\partopsep}{0pt}}%
  \fontsize{9pt}{11.5pt}\bodyfontlight\color{text}\raggedright}%
  {\end{list}\vspace{1mm}}
\newcommand{\cvitem}[2]{\item[{\fontsize{8.5pt}{10pt}\bodyfont\color{graytext}#1\hfill}]#2}
\newcommand{\skilllabel}[1]{{\bodyfont\color{darktext}#1}}
\renewcommand{\textbf}[1]{{\bodyfont\bfseries\color{darktext}#1}}


\photo[circle,edge,left]{Tobias_Staal_square}
\name{Tobias}{Stål}
\position{Geoscientist and engineer}
\address{Tasmania, Australia}
\mobile{+61 4 4749 4089}
\email{tobias.staal@utas.edu.au}
\homepage{tobbetripitaka.github.io/Tobias\_Staal}
\github{TobbeTripitaka}
\makecvfooter{8 October 2026}{Tobias Stål · Curriculum Vitae}{\thepage}

\begin{document}
\makecvheader

\vspace{2mm}
{\fontsize{9.5pt}{12.5pt}\bodyfontlight\color{text}Geophysicist working on Antarctica's ice-covered interior. I integrate geology, geophysics and statistical modelling to understand interactions between the ice sheet and the solid Earth, with a focus on geothermal heat, subglacial geology and instrumentation. Background in geology and geomorphology from the Arctic and Europe, and an earlier professional career in sound, lighting and electrical engineering.\par}

\cvsection{Academic positions}
\begin{cvlist}
\cvitem{2022 – Present}{\textbf{Research Associate in Computational Physics (Geophysics)}, ACEAS, University of Tasmania, Tasmania, Australia}
\cvitem{2021 – 2022}{\textbf{Research Fellow in Antarctic Seismology}, University of Tasmania, Tasmania, Australia}
\cvitem{2020 – 2021}{\textbf{Research Assistant}, University of Tasmania, Tasmania, Australia}
\cvitem{2016 – 2021}{\textbf{Casual tutor, demonstrator and research assistant}, University of Tasmania, Tasmania, Australia}
\end{cvlist}

\cvsection{Education}
\begin{cvlist}
\cvitem{2021}{\textbf{PhD, Geophysics, Geology and Computational Methods}, University of Tasmania, Australia\par{\color{graytext}Thesis: \textit{The Antarctic lithosphere revealed by multivariate analysis} \newline Supervisors: AM Reading, JA Halpin, JM Whittaker \newline Coursework in statistics and computational geophysics}}
\cvitem{2015}{\textbf{MSc, Geophysics}, University of Copenhagen, Denmark\par{\color{graytext}Thesis: \textit{High-resolution S- and P-wave seismic mapping of near-surface chalk deposits} \newline Supervisor: L Nielsen \newline Coursework in geophysics, stochastic modelling and sedimentary basins}}
\cvitem{2011}{\textbf{BSc, Geology}, University of Copenhagen, Denmark\par{\color{graytext}Including studies at UNIS, Svalbard \newline Thesis: \textit{Processes and development of the Faroe North Slide Complex} \newline Supervisors: GK Pedersen, T Nielsen \newline Coursework in geology, palaeontology, geomorphology, Quaternary geology and Arctic marine geology}}
\cvitem{1998}{\textbf{AD, Sound Engineering}, Piteå Music Conservatory (Luleå University of Technology), Piteå, Sweden\par{\color{graytext}Coursework in electronics, acoustics, radio broadcasting, music theory and recording technology \newline Including exchange studies at VGIK, Moscow, Russia}}
\cvitem{Ongoing}{\textbf{Unincorporated studies (162.5 ECTS)}\par{\color{graytext}Linguistics (linguistic typology, sociolinguistics, African languages, Nordic studies, pragmatics, grammar) \newline Geology: climate change, non-Nordic geology (the American Cordillera), Arctic petroleum provinces, geohazards, GIS and remote sensing}}
\end{cvlist}

\cvsection{Peer-reviewed publications}
\begin{cvlist}
\cvitem{2026}{JA Halpin, NR Daczko, JA Mulder, A Maritati, \textbf{T Stål} — Ediacaran–Cambrian High‐Strain Melt‐Present Deformation of the Archean Charcot Province, East Antarctica. \textit{Tectonics}, 45(8), e2026TC009479. \href{https://doi.org/10.1029/2026TC009479}{doi:10.1029/2026TC009479}}
\cvitem{2026}{MC Manassero, K Selway, \textbf{T Stål}, M Scheiter, M Loesing, F McCormack, JA Halpin, B Kulessa, AM Reading — Bed Topography and Subglacial Conditions of Denman Glacier, East Antarctica: Insights From Magnetotelluric Data and Interdisciplinary Studies. \textit{Journal of Geophysical Research: Earth Surface}, 131(7), e2026JF009277. \href{https://doi.org/10.1029/2026JF009277}{doi:10.1029/2026JF009277}}
\cvitem{2026}{ID Kelly, AM Reading, \textbf{T Stål}, B Kulessa, A García‐Jerez, J Piña‐Flores, E Paolucci, A Tanzini, RJ Turner, JC Magyar et al. — Determining the Character of Subglacial Sediments in the Ice‐Bedrock Interface Zone of Antarctica Using Horizontal‐to‐Vertical Spectral Ratios (HVSRs) of Seismic Ambient Noise. \textit{Journal of Geophysical Research: Solid Earth}, 131(6), e2025JB033029. \href{https://doi.org/10.1029/2025JB033029}{doi:10.1029/2025JB033029}}
\cvitem{2026}{FS McCormack, \textbf{T Stål}, N Shao, E MacKie, A Fabela Hinojosa, M Lösing, J Roberts, S Ehrenfeucht, C Dow — Synthetic bed topographies for Antarctica and their utility in ice sheet modelling. \textit{Philosophical Transactions of the Royal Society A: Mathematical, Physical and Engineering Sciences}, 384(2319), 20240537. \href{https://doi.org/10.1098/rsta.2024.0537}{doi:10.1098/rsta.2024.0537}}
\cvitem{2026}{D Arora, NC Pant, \textbf{T Stål}, A Naik, M Pandey, R Gupta — Geothermal heat flow in the Princess Elizabeth Land sector, East Antarctica. \textit{Earth-Science Reviews}, 277, 105475. \href{https://doi.org/10.1016/j.earscirev.2026.105475}{doi:10.1016/j.earscirev.2026.105475}}
\cvitem{2026}{M Al-Aghbary, M Sobh, \textbf{T Stål}, MO Awaleh, M Jalludin, C Gerhards — Uncertainty quantification using clustering-based quantile regression forests: with an application to improving geothermal heat flow predictions. \textit{Geophysical Journal International}, 245(2), ggag113. \href{https://doi.org/10.1093/gji/ggag113}{doi:10.1093/gji/ggag113}}
\cvitem{2026}{M Lösing, W Colgan, \textbf{T Stål}, J Ebbing, AG Busck, T Zhang, HL Seroussi, F McCormack, D Fahrner, L Stearns et al. — Community heat flow recommendations: suitable basal boundary conditions for Greenland and Antarctica in ISMIP7. \textit{GEUS Bulletin}, 62. \href{https://doi.org/10.34194/r0w9rf81}{doi:10.34194/r0w9rf81}}
\cvitem{2025}{M Lösing, A Aitken, J Ebbing, JA Halpin, L Li, M Moorkamp, A Reading, \textbf{T Stål} — Linking Tectonics and Crustal Thermal Properties in Southwestern Australia and East Antarctica Through Coupled Gravity and Magnetic Analysis. \textit{Journal of Geophysical Research: Solid Earth}, 130(12), e2024JB030770. \href{https://doi.org/10.1029/2024JB030770}{doi:10.1029/2024JB030770}}
\cvitem{2025}{NJ Abram, A Purich, MH England, FS McCormack, JM Strugnell, DM Bergstrom, TR Vance, \textbf{T Stål}, B Wienecke, P Heil et al. — Emerging evidence of abrupt changes in the Antarctic environment. \textit{Nature}, 644(8077), 621–633. \href{https://doi.org/10.1038/s41586-025-09349-5}{doi:10.1038/s41586-025-09349-5}}
\cvitem{2024}{\textbf{T Stål}, JA Halpin, JW Goodge, AM Reading — Geology Matters for Antarctic Geothermal Heat. \textit{Geophysical Research Letters}, 51(13), e2024GL110098. \href{https://doi.org/10.1029/2024GL110098}{doi:10.1029/2024GL110098}}
\cvitem{2022}{FS McCormack, JL Roberts, CF Dow, \textbf{T Stål}, JA Halpin, AM Reading, MJ Siegert — Fine‐Scale Geothermal Heat Flow in Antarctica Can Increase Simulated Subglacial Melt Estimates. \textit{Geophysical Research Letters}, 49(15), e2022GL098539. \href{https://doi.org/10.1029/2022GL098539}{doi:10.1029/2022GL098539}}
\cvitem{2022}{\textbf{T Stål}, AM Reading, S Fuchs, JA Halpin, M Lösing, RJ Turner — Properties and biases of the global heat flow compilation. \textit{Frontiers in Earth Science}, 10, 963525. \href{https://doi.org/10.3389/feart.2022.963525}{doi:10.3389/feart.2022.963525}}
\cvitem{2022}{AM Reading, \textbf{T Stål}, JA Halpin, M Lösing, J Ebbing, W Shen, FS McCormack, CS Siddoway, D Hasterok — Antarctic geothermal heat flow and its implications for tectonics and ice sheets. \textit{Nature Reviews Earth \& Environment}, 3(12), 814–831. \href{https://doi.org/10.1038/s43017-022-00348-y}{doi:10.1038/s43017-022-00348-y}}
\cvitem{2022}{PE Morse, AM Reading, \textbf{T Stål} — Exploratory Volumetric Deep Earth Visualization by 2.5D Interactive Compositing. \textit{IEEE Transactions on Visualization and Computer Graphics}, 28(7), 2641–2653. \href{https://doi.org/10.1109/TVCG.2020.3037226}{doi:10.1109/TVCG.2020.3037226}}
\cvitem{2021}{\textbf{T Stål}, AM Reading, JA Halpin, JM Whittaker — Antarctic Geothermal Heat Flow Model: Aq1. \textit{Geochemistry, Geophysics, Geosystems}, 22(2), e2020GC009428. \href{https://doi.org/10.1029/2020GC009428}{doi:10.1029/2020GC009428}}
\cvitem{2021}{G Sanchez, JA Halpin, M Gard, D Hasterok, \textbf{T Stål}, T Raimondo, S Peters, A Burton‐Johnson — PetroChron Antarctica: A Geological Database for Interdisciplinary Use. \textit{Geochemistry, Geophysics, Geosystems}, 22(12), e2021GC010154. \href{https://doi.org/10.1029/2021GC010154}{doi:10.1029/2021GC010154}}
\cvitem{2020}{\textbf{T Stål}, AM Reading, JA Halpin, SJ Phipps, JM Whittaker — The Antarctic Crust and Upper Mantle: A Flexible 3D Model and Software Framework for Interdisciplinary Research. \textit{Frontiers in Earth Science}, 8, 577502. \href{https://doi.org/10.3389/feart.2020.577502}{doi:10.3389/feart.2020.577502}}
\cvitem{2020}{\textbf{T Stål}, AM Reading — A Grid for Multidimensional and Multivariate Spatial Representation and
                        Data Processing. \textit{Journal of Open Research Software}, 8(1), 2. \href{https://doi.org/10.5334/jors.287}{doi:10.5334/jors.287}}
\cvitem{2019}{\textbf{T Stål}, AM Reading, JA Halpin, JM Whittaker — A Multivariate Approach for Mapping Lithospheric Domain Boundaries in East Antarctica. \textit{Geophysical Research Letters}, 46(17-18), 10404–10416. \href{https://doi.org/10.1029/2019GL083453}{doi:10.1029/2019GL083453}}
\cvitem{2019}{PE Morse, AM Reading, \textbf{T Stål} — Well-Posed Geoscientific Visualization Through Interactive Color Mapping. \textit{Frontiers in Earth Science}, 7, 274. \href{https://doi.org/10.3389/feart.2019.00274}{doi:10.3389/feart.2019.00274}}
\end{cvlist}

\cvsection{Submitted and in review}
\begin{cvlist}
\cvitem{2026}{MC Manassero, K Selway, AM Reading, N Askey-Doran, JR Peacock, F Ramirez, \textbf{T Stål} — Magnetotellurics, mantle viscosity, integrated geophysics, glacial isostatic adjustment, mantle water content, East Antarctica. \textit{Journal of Geophysical Research: Solid Earth}. In review.}
\end{cvlist}

\cvsection{Outreach}
\begin{cvlist}
\cvitem{2024}{\textbf{T Stål}, FS McCormack, AM Reading, N Askey-Doran, JA Halpin, M Lösing — Geothermal Heat Shapes the Antarctic Ice Sheet From Below. \textit{Frontiers for Young Minds}, 11, 1178537. \href{https://doi.org/10.3389/frym.2023.1178537}{doi:10.3389/frym.2023.1178537}}
\end{cvlist}

\cvsection{Datasets}
\begin{cvlist}
\cvitem{2024}{Global Heat Flow Data Assessment Group, S Fuchs, F Neumann, B Norden, E Balkan-Pazvantoglu, S Elbarbary, A Petrunin, G Beardsmore, R Harris, …, \textbf{T Stål} — The Global Heat Flow Database: Release 2024 [Dataset]. \textit{GFZ Data Services}. \href{https://doi.org/10.5880/fidgeo.2024.014}{doi:10.5880/fidgeo.2024.014}}
\cvitem{2024}{AM Reading, \textbf{T Stål}, I Kelly — Ice Sheets and Interactions, Dronning Maud Land and Coats Land [Dataset]. \textit{International Federation of Digital Seismograph Networks}. \href{https://doi.org/10.7914/J9CQ-XD60}{doi:10.7914/J9CQ-XD60}}
\cvitem{2023}{Global Heat Flow Data Assessment Group, S Fuchs, F Neumann, B Norden, G Beardsmore, P Chiozzi, W Colgan, AP Anguiano Dominguez, MRA Duque, …, \textbf{T Stål} et al. — The Global Heat Flow Database: Update 2023 [Dataset]. \textit{GFZ Data Services}. \href{https://doi.org/10.5880/fidgeo.2023.008}{doi:10.5880/fidgeo.2023.008}}
\cvitem{2023}{\textbf{T Stål}, AM Reading, S Kupis — Ice Sheets and Interactions, East Antarctica [Dataset]. \textit{International Federation of Digital Seismograph Networks}. \href{https://doi.org/10.7914/74EN-YT90}{doi:10.7914/74EN-YT90}}
\cvitem{2021}{\textbf{T Stål}, AM Reading, S Kupis, J Newlands — Casey-Wilkins Adaptable Array [Dataset]. \textit{International Federation of Digital Seismograph Networks}. \href{https://doi.org/10.7914/SN/Z9_2021}{doi:10.7914/SN/Z9\_2021}}
\cvitem{2020}{\textbf{T Stål}, AM Reading, JA Halpin, SJ Phipps, J Whittaker — The Antarctic crust and upper mantle: a flexible 3D model and framework for interdisciplinary research [Dataset]. \textit{PANGAEA}. \href{https://doi.org/10.1594/PANGAEA.918549}{doi:10.1594/PANGAEA.918549}}
\cvitem{2015}{AM Reading, \textbf{T Stål} — CAD Seismic Deployment, Casey-Davis Region, East Antarctica [Dataset]. \textit{International Federation of Digital Seismograph Networks}. \href{https://doi.org/10.7914/NRW9-9825}{doi:10.7914/NRW9-9825}}
\end{cvlist}

\cvsection{Software and code}
\begin{cvlist}
\cvitem{2022}{\textbf{T Stål} — Fast Calculation of Ripley's K and L Functions on a Sphere [Software]. \textit{Zenodo}. \href{https://doi.org/10.5281/zenodo.6865005}{doi:10.5281/zenodo.6865005}}
\cvitem{2019}{\textbf{T Stål} — agrid: a grid for multidimensional and multivariate spatial modelling [Software]. \textit{Zenodo}. \href{https://doi.org/10.5281/zenodo.2553966}{doi:10.5281/zenodo.2553966}}
\end{cvlist}

\cvsection{Selected conferences, workshops and hackathons}
\begin{cvlist}
\cvitem{2026}{EGU General Assembly, Vienna, Austria\par{\color{graytext}– Aq2 and Kq2, refined geothermal heat flow models from multivariate observables for application to ice sheet modelling}}
\cvitem{2025}{geoBRIDGE Working Group Hackathon, Perth, Australia\par{\color{graytext}– Participant and project contributor}}
\cvitem{2025}{EGU General Assembly, Vienna, Austria\par{\color{graytext}– Seismicity of Denman Glacier: constraints on geometry and dynamics}}
\cvitem{2024}{SCAR Open Science Conference, Pucón, Chile\par{\color{graytext}– Aq2: a refined geothermal heat flow map of Antarctica from multivariate observables for application to ice sheet modelling \newline – Exploring Atlas Antarktiki: a distinctive repository of historical Antarctic and Southern Ocean data \newline – Convened three scientific sessions on Antarctic geology and geophysics, and co-organised a side meeting to organise collaborations at Bunger Hills}}
\cvitem{2023}{EGU General Assembly, Vienna, Austria\par{\color{graytext}– Using information entropy to optimise and communicate certainty of continental scale tectonic models}}
\cvitem{2023}{Maydena Subglacial Hydrology and Geology Workshop, Maydena, Australia\par{\color{graytext}– The subglacial heat budget – a conceptual investigation \newline – Organised and co-chaired the workshop}}
\cvitem{2021}{AGU Fall Meeting, New Orleans and online\par{\color{graytext}– Modelling Antarctic subglacial geothermal heat transfer in higher resolution, with reduced uncertainty and assumptions addressed \newline – Co-convened a session on Antarctic geothermal heat}}
\cvitem{2021}{Australian Earth Sciences Convention (AESC) "Core to Cosmos", online\par{\color{graytext}– Data-driven tectonic regionalisation of Antarctica: appreciate the similarity \newline – Convened two sessions on data science and machine learning}}
\cvitem{2020}{SCAR Open Science Conference: Antarctic – Global Connections, online\par{\color{graytext}– Towards a reconciled heat flow map for Antarctica: Aq1 \newline – The Wilkes Land sector including the Aurora Basin, and its most probable subglacial geology}}
\cvitem{2020}{AGU Fall Meeting, online\par{\color{graytext}– Multivariate geothermal heat flow models of Antarctica, towards Aq2}}
\cvitem{2019}{Geological Society of Australia SGTSG Meeting, Port Lincoln, Australia\par{\color{graytext}– The probable geology of the Aurora Basin, East Antarctica}}
\cvitem{2019}{XIII International Symposium on Antarctic Earth Sciences (ISAES), Incheon, Korea\par{\color{graytext}– Linking Antarctic geological observations and geophysical data in a probabilistic space}}
\cvitem{2018}{POLAR2018 Open Science Conference, Davos, Switzerland\par{\color{graytext}– A multi-domain lithospheric model of East Antarctica}}
\cvitem{2018}{SCAR TACtical heat flow workshop, Hobart, Australia\par{\color{graytext}– Talk and poster on the tectonic structure of Antarctica}}
\cvitem{2018}{IEEE Antarctic and Southern Ocean Forum (ASOF), Hobart, Australia\par{\color{graytext}– A frame for processing multivariate Antarctic solid Earth data}}
\cvitem{2018}{DataTas seminar, Hobart, Australia\par{\color{graytext}– Software construction tools make everything so easy!}}
\cvitem{2018}{Agile* Subsurface Hackathon, Copenhagen, Denmark\par{\color{graytext}– Winner of the People's Choice Award}}
\cvitem{2017}{AGU Fall Meeting, New Orleans, USA\par{\color{graytext}– Combined constraints on the structure and physical properties of the East Antarctic lithosphere from geology and geophysics}}
\cvitem{2015}{Workshop on 3D Reflection Seismic Processing Using Open-Source Software, Rice University, Houston, Texas, USA}
\end{cvlist}

\cvsection{Awards and recognition}
\begin{cvlist}
\cvitem{2023}{Royal Society of Tasmania Doctoral Award}
\cvitem{2021}{Wiley Top Cited Article}
\cvitem{2020}{PhD Award for Outstanding Performance during Postgraduate Studies}
\cvitem{2019}{Progress in Inclusion, Diversity \& Equity}
\end{cvlist}

\cvsection{Memberships and service}
\begin{cvlist}
\cvitem{2023 – Present}{Member, International Heat Flow Commission (IHFC), IUGG}
\cvitem{2021 – Present}{Co-chair, SCAR INSTANT Heat Flow Subcommittee}
\cvitem{2020}{Organising committee, Australian Antarctic Research Conference (AARC)}
\cvitem{2019}{Organising committee, Geological Society of Australia SGTSG Meeting}
\cvitem{2017 – 2020}{Member, School of Natural Sciences Inclusion, Diversity \& Equity Committee}
\end{cvlist}

\cvsection{Selected non-academic work experience}
\begin{cvlist}
\cvitem{2021 – Present}{\textbf{GRIT Facility Manager – Instrumentation}, University of Tasmania, Tasmania, Australia\par{\color{graytext}Instrumentation management and Antarctic fieldwork support}}
\cvitem{2019 – 2021}{\textbf{Geoscientist}, Geoneon, Australia\par{\color{graytext}Marine geoscience, GIS and data compilation}}
\cvitem{2011 – 2012}{\textbf{Invited technician and video designer}, Al-Harah Theatre, Palestine}
\cvitem{2010 – 2015}{\textbf{Field Leader and Geophysicist}, COWI, Denmark\par{\color{graytext}MEP, GPR and TEM fieldwork, and GIS-based survey design and interpretation}}
\cvitem{2010}{\textbf{Field Scientist}, Norwegian Polar Institute, Svalbard\par{\color{graytext}Supporting biological fieldwork}}
\cvitem{2009 – 2016}{\textbf{Lighting Designer}, Rapid Eye, Denmark\par{\color{graytext}Including the award-winning performance \textit{QUIPROQUO}}}
\cvitem{2009 – 2014}{\textbf{Lighting Designer}, STÜ, Estonia\par{\color{graytext}Including touring}}
\cvitem{2008 – 2009}{\textbf{Technical Producer}, National Theatre of Ghana, Ghana}
\cvitem{2005 – 2007}{\textbf{Sound/Video/Light Designer and Technician}, The Swedish National Touring Theatre, Sweden\par{\color{graytext}And shorter projects until 2011}}
\cvitem{2004 – 2006}{\textbf{Sound Engineer}, Circus Cirkör, Sweden and world tour}
\cvitem{2003 – 2004}{\textbf{Sound Designer}, The Royal Dramatic Theatre, Sweden}
\cvitem{1998 – 2003}{\textbf{Stage Manager}, Uppsala Stadsteater, Sweden}
\cvitem{1997 – 2004}{\textbf{Sound Engineer}, Swedish Radio (P1, P3), Sweden}
\cvitem{1993 – 1996}{\textbf{Plumbing Apprentice}, Linds Rör AB, Sweden}
\end{cvlist}

\cvsection{Art collaborations (as a scientist)}
\begin{cvlist}
\cvitem{2025}{T Stål, AM Reading, K Verell — \textit{Seismic Dance Party}, Beaker Street (Hobart). Collaborative science–art performance with AuScope, sound designer K Verell and DJs Huh Pink, translating Denman Glacier seismic data into an immersive audio-visual dancefloor experience.}
\cvitem{2025}{L Bleach — \textit{An Infrasonorous Archipelago}, Walk\&Talk Biennial "Gestures of Abundance" (Azores). Using seismo-volcanic infrasound from the Azores to create tactile and sonic works that tune audiences to the Earth’s low-frequency vibrations.}
\cvitem{2025}{M Nieboer — \textit{Through Our Eyes}, Australian Embassy (Paris). Exhibition on women’s experiences and narratives, the culmination of two years of work; I contributed cartographic design and geospatial visualisation.}
\cvitem{2023}{L Bleach — \textit{soft rock monitoring (the slow seismogenic zone)}, EARTHEN. Better Nature, Cement Fondu (Sydney). Collaboration with Lucy Bleach and Geoscience Australia, transmitting vibrational signals from the SYDS soft-rock urban seismic monitoring site into the gallery architecture via transducer speakers.}
\cvitem{2021}{L Bleach — \textit{Attenuated ground (the slow seismogenic zone)}, TarraWarra Biennial 2021: Slow Moving Waters (TarraWarra Museum of Art, Victoria). Integrating double bass, cast toffee, seismometer, tactile transducers and slow-earthquake signals from the Japan Trench to render seismogenic motion as resonant sound and vibration.}
\end{cvlist}

\cvsection{Other skills and experience}
\begin{cvlist}
\cvitem{\skilllabel{Coding}}{Python (scikit-learn, OpenCV, TensorFlow, dask, xarray, geopandas), GIS (QGIS, ArcGIS, GDAL/OGR), TeX/LaTeX (TikZ), Bash, Linux, macOS, git, SCons, Madagascar, Seismic Unix, SQL and MAX. Some HTML, C and JavaScript (learning). Microsoft Office, Adobe Creative Suite, Blender (BlenderGIS), GIMP, Inkscape, CAD (Fusion). Working but outdated knowledge of MATLAB, SAC, Emacs Lisp and BASIC.}
\cvitem{\skilllabel{Certificates}}{Coastal Skipper Certificate (SWE); Marine HF/VHF Radio Certificate (SWE); firearms licence and hunting certificate (SWE); PADI Advanced Open Water Diver (including dry suit and Nitrox); Wilderness First Aid (AUS); Australian Approved Arrangements Accreditation Biosecurity (AUS); trip leader and terrain-vehicle driver for the Australian Antarctic Division (e.g. Casey 2023)}
\cvitem{\skilllabel{Languages}}{English (C2), Swedish (native), Danish (fluent), French (intermediate), Russian (functional), Asante Twi (functional), Norwegian (reading and understanding), Esperanto (inactive), Spanish (learning), Yiddish (learning), Estonian (learning)}
\cvitem{\skilllabel{Online courses}}{Drawing Nature, Science and Culture: Natural History Illustration NHI101x (University of Newcastle, 2020); 3D Earth spring school (Kiel University, 2021); Presenting Data and Information by Edward Tufte (2025)}
\cvitem{\skilllabel{Practical skills}}{Welding, woodwork, diesel engines, sewing, cooking and baking (enquire about my chocolate mousse)}
\cvitem{\skilllabel{Other interests}}{Music, literature and theatre. I also love old boats, particularly tugs.}
\end{cvlist}


\end{document}
